%! AP Statistics — Unit 4, Lesson 4.2 — Warm-Up KEY
\documentclass[10pt]{article}
\usepackage{apstats-article}
\usepackage{apstats-key}

\begin{document}

\pageheader{Unit 4, Lesson 4.2}{Warm-Up --- KEY}
\namedateperiod

\vspace{0.3in}

{\large\bfseries\color{navy} Part 1 \quad Spiral Review}

\vspace{0.12in}

\begin{enumerate}

    \item From Lesson 4.1 --- for each pattern, write the statistical question it suggests
    and state whether chance is a plausible explanation.

    \vspace{8pt}
    \renewcommand{\arraystretch}{1.8}
    \begin{tabularx}{\linewidth}{@{} X p{4cm} p{2.5cm} @{}}
    \textbf{Pattern} & \textbf{Statistical question} & \textbf{Chance plausible?} \\
    \hline
    A batter gets 6 hits in his last 7 at-bats (career average .280). &
    \ans{Is this batter's hit rate during this stretch higher than his career average?} & \ans{Yes, but unlikely} \\
    A new drug trial shows 18 of 20 patients improved. &
    \ans{Does this drug improve patient outcomes at a rate higher than chance?} & \ans{Possibly not} \\
    \end{tabularx}

    \vspace{0.18in}

    \item
    \begin{enumerate}[label=(\alph*), itemsep=6pt]
        \item \ans{A: 8/40 = 0.20; \quad B: 12/40 = 0.30; \quad C: 11/40 = 0.275; \quad D: 9/40 = 0.225}
        \item \ans{0.25 for each section (1/4).}
        \item \ans{The results are close to 0.25 for all sections. The differences are small and likely due to random variation with only 40 spins. The data do not strongly suggest the spinner is unfair.}
    \end{enumerate}

\end{enumerate}

\vspace{0.35in}

{\large\bfseries\color{navy} Part 2 \quad Looking Ahead}

\vspace{0.12in}

\noindent Today we build a tool to estimate probabilities before we have formulas.

\vspace{0.14in}

\begin{tcolorbox}[colback=greenbg, colframe=greenacc, boxrule=0.8pt, arc=2mm,
    left=4mm, right=4mm, top=3mm, bottom=3mm]
    \small
    \begin{itemize}[leftmargin=1.3em, itemsep=8pt]
  \item A \textbf{simulation} uses a chance device to model a random process.
  \item We estimate probability by computing the \textbf{relative frequency} of the outcome over many trials.
  \item More trials $\to$ more accurate estimate (Law of Large Numbers).
    \end{itemize}
\end{tcolorbox}

\vspace{0.2in}

\noindent\textcolor{slate}{\small If you wanted to estimate the probability that a thumbtack
lands point-up when dropped, how would you go about it?}\\[8pt]
\begin{teachernote}
Accept reasonable plans. Strong responses will include: drop the thumbtack many times, record how many times it lands point-up, divide by total drops to get relative frequency. Bonus if student mentions needing many trials for an accurate estimate.
\end{teachernote}

\end{document}
